\section{Conclusion and Future Work}
\label{sec:conclusion}

The study specifies a full path from the public Tetouan series to an
operating alert. Ten-minute zone measurements are summed and aggregated to
hourly city demand. Calendar, lag, rolling and weather features are built
without using the future. A seasonal-naive baseline, ridge regression, a
random forest, XGBoost and LightGBM are compared on a later test block that
was not used for training or for model selection. Beside MAE, RMSE and WAPE,
the report records peak magnitude, peak timing and breach anticipation, which
are the scores that speak to reserve-margin decisions.

Once Table~\ref{tab:scores} and Table~\ref{tab:peak} are filled, the
conclusion should state the best model by MAE and whether that same model
also leads on breach anticipation. The operating recommendation follows the
hour profile in Figure~\ref{fig:importance}: arm reserve one to two hours
before the clock window that most often exceeds \(\tau\), with the evening
hours from 18:00 to 21:00 as the window this design is built to protect.

Two extensions stay outside the present pipeline. The first is a quantile
forecast, so the alert can be raised from an upper quantile rather than from
a single point prediction. The second is a multi-horizon forecast at 1 hour,
6 hours and 24 hours, using a separate model or a direct horizon feature for
each lead time. Both extensions reuse the same chronological split and the
same peak metrics.
